% from ega3-1-fr.tex
\oldpage[0\textsubscript{III}]{45}
La proposition (\hyperref[0.11.8.5-fr]{11.8.5}) ramène le calcul de la
cohomologie de $\mathrm L^\bullet(A,B;\Gamma)$ à celui de la cohomologie de
$\mathrm C^\bullet(A,B;\Gamma)$. Rappelons d'autre part que cette dernière
est, en vertu du th. d'Eilenberg-Zilber (G, I, 3.10.2), canoniquement
isomorphe à la cohomologie du complexe de cochaînes défini comme suit : on
forme le complexe de chaînes $\mathrm P_\bullet(A,B)$, constitué par les
combinaisons linéaires des $(\sigma,\tau)\in\Sigma(A)\times\Sigma(B)$ tels
que $\sigma$ et $\tau$ aient le même degré; la différentielle de ce complexe
est donnée par
\[
d(\sigma,\tau)=\sum_{j,k}(-1)^{j+k}(\sigma_j,\tau_k)
\]
si
\[
d\sigma=\sum_j(-1)^j\sigma_j,\qquad
d\tau=\sum_k(-1)^k\tau_k;
\]
on a alors deux homomorphismes canoniques de complexes
\[
f:\mathrm P_\bullet(A,B)\longrightarrow
\mathrm C_\bullet(A)\otimes\mathrm C_\bullet(B),\qquad
g:\mathrm C_\bullet(A)\otimes\mathrm C_\bullet(B)
\longrightarrow\mathrm P_\bullet(A,B),
\]
et on démontre (\emph{loc. cit.}) qu'il y a des homotopies $h$, $h'$ telles
que
\[
f\circ g-1=d\circ h+h\circ d,\qquad
g\circ f-1=d\circ h'+h'\circ d.
\]
En outre, on a
\[
f(\sigma,\tau)\in\mathrm C_\bullet(|\sigma|)
\otimes\mathrm C_\bullet(|\tau|),\qquad
g(\sigma,\tau)\in\mathrm P_\bullet(|\sigma|,|\tau|)
\]
et les homotopies $h$, $h'$ peuvent être prises telles que
\[
h(\sigma,\tau)\in\mathrm C_\bullet(|\sigma|)
\otimes\mathrm C_\bullet(|\tau|),\qquad
h'(\sigma,\tau)\in\mathrm P_\bullet(|\sigma|,|\tau|).
\]
Ce point provient de ce que la définition de $h(\sigma,\tau)$ et
$h'(\sigma,\tau)$ peut se faire par récurrence sur la somme des degrés de
$\sigma$ et $\tau$, et de ce que les
\[
H^q(\mathrm C_\bullet(|\sigma|)\otimes\mathrm C_\bullet(|\tau|))
\quad\text{et}\quad
H^q(\mathrm P_\bullet(|\sigma|,|\tau|))
\]
sont nuls pour $q>0$ (\emph{loc. cit.}); on raisonne alors comme dans
(\hyperref[0.11.8.3-fr]{11.8.3}) et la conclusion en résulte.

On définit alors $\mathrm P^\bullet(A,B;\Gamma)$ comme l'ensemble des
familles $\lambda=(\lambda(\sigma,\tau))$ où $(\sigma,\tau)$ parcourt les
couples dont les termes ont même degré, avec
$\lambda(\sigma,\tau)\in\Gamma_{\sigma,\tau}$, et comme on a
\[
d\sigma=\sum_j(-1)^j\sigma_j,\qquad
d\tau=\sum_k(-1)^k\tau_k,
\]
\[
d\lambda(\sigma,\tau)=
\sum_{j,k}(-1)^{j+k}\lambda(\sigma_j,\tau_k)
\]
est défini et donne la différentielle du complexe
$\mathrm P^\bullet(A,B;\Gamma)$. Cela étant, les applications
\[
{}'f:\lambda\mapsto\lambda\circ f,\quad
{}'g:\lambda\mapsto\lambda\circ g,\quad
{}'h:\lambda\mapsto\lambda\circ h,\quad
{}'h':\lambda\mapsto\lambda\circ h'
\]
sont toutes définies en vertu des remarques qui précèdent;
$\mathop{}'f\circ{}'g$ et $\mathop{}'g\circ{}'f$ sont donc homotopes à
l'identité, d'où l'isomorphisme cherché entre la cohomologie de
$\mathrm C^\bullet(A,B;\Gamma)$ et celle de
$\mathrm P^\bullet(A,B;\Gamma)$.
\end{env}

\begin{env}[11.8.7]
\label{0.11.8.7-fr}
\emph{Remarque}. Le même raisonnement que dans
(\hyperref[0.11.8.3-fr]{11.8.3}), mais appliqué à
$\mathrm C_\bullet(A)$ et $\mathrm L_\bullet(A)$, montre que si $j$ et $p$
sont définis comme dans (\hyperref[0.11.8.1-fr]{11.8.1}), $f=j\circ p$
vérifie encore une relation (\hyperref[0.11.8.3.1-fr]{11.8.3.1}), avec
$|h(\sigma)|\subset|\sigma|$, d'où on déduit comme dans
(\hyperref[0.11.8.5-fr]{11.8.5}) un isomorphisme de la cohomologie de
$\mathrm L^\bullet(A;\Gamma)$ sur celle de
$\mathrm C^\bullet(A;\Gamma)$, ces deux complexes étant définis de façon
évidente. C'est le résultat dont la démonstration est esquissée dans
(G, I, 3.8.1).
\end{env}

\begin{env}[11.8.8]
\label{0.11.8.8-fr}
Reprenons maintenant les notations et hypothèses de
(\hyperref[0.11.8.2-fr]{11.8.2}), et considérons un complexe
$\Gamma^\bullet=(\Gamma^k)$ de systèmes de coefficients sur
$\Sigma(A)\times\Sigma(B)$ : pour chaque $(\sigma,\tau)$, les
$\Gamma^k_{\sigma,\tau}$ forment donc un complexe de groupes abéliens
$(k\in\mathbb Z)$, et on a les diagrammes commutatifs
\[
\xymatrix{
\Gamma^k_{\sigma,\tau}\ar[r]\ar[d] &
\Gamma^{k+1}_{\sigma,\tau}\ar[d]\\
\Gamma^k_{\sigma',\tau'}\ar[r] &
\Gamma^{k+1}_{\sigma',\tau'}
}
\]
